\documentclass[12pt]{article}

\usepackage[margin=1.25in]{geometry}
\usepackage{setspace}
\usepackage{hyperref}
\usepackage{booktabs}
\usepackage{footnote}
\usepackage{longtable}

\hypersetup{
    colorlinks=true,
    linkcolor=blue,
    citecolor=blue,
    urlcolor=blue
}

\doublespacing

\title{British Precedents for Union Civil War Finance:\\
Pitt's Fiscal--Military State and Lincoln's Treasury}
\author{Notes prepared for the fiscal history project}
\date{\today}

\begin{document}

\maketitle
\thispagestyle{empty}

\newpage

\tableofcontents

\newpage

%-----------------------------------------------------------------------
\section{Introduction}
%-----------------------------------------------------------------------

The Union's fiscal response to the Civil War---legal tender notes,
long-term funded bonds, a new income tax, excise duties, a rationalized
banking system, and a sinking fund backed by dedicated customs
revenue---had no precedent in prior American experience. But nearly every
element had a close precedent in British war finance during the French
Revolutionary and Napoleonic Wars (1793--1815). Whether Union
policymakers drew on these British models consciously, absorbed them
through the intellectual atmosphere of Anglo-American political economy,
or simply arrived at similar solutions under similar pressures, the
structural parallels are striking and, in several cases, unmistakable.

This document traces seven major parallels between British fiscal policy
during 1793--1815 and Union fiscal policy during 1861--1865, and
identifies the channels---intellectual, institutional, and
structural---through which British experience shaped American choices.
The parallels run deeper than surface resemblance: they share a common
logic rooted in the fiscal demands of existential warfare and the
political economy of credible commitment to debt service.

%-----------------------------------------------------------------------
\section{The Fiscal--Military State as Template}
%-----------------------------------------------------------------------

The concept of the ``fiscal--military state'' was developed by John
Brewer (\textit{The Sinews of Power}, 1989) to describe England's
transformation after the Glorious Revolution of 1688. Its defining
features were:

\begin{enumerate}
\item A permanent funded national debt, managed by the Treasury and
  serviced by dedicated tax revenues.
\item A central bank (the Bank of England, founded 1694) that served as
  fiscal agent, government banker, and manager of the public debt.
\item A professional revenue bureaucracy---especially the Customs and
  Excise services---capable of raising large sums reliably.
\item A tax system that combined customs duties, excise taxes on
  domestic manufactures and consumption, and (after 1799) a direct
  income tax.
\item A sinking fund, established by Walpole in 1717 and reconstituted
  by Pitt in 1786, designed to retire debt systematically from
  earmarked revenues.
\end{enumerate}

Patrick O'Brien (\textit{Power with Profit}, 1998; and numerous
articles) documented how this fiscal--military state enabled Britain to
sustain debt-to-GDP ratios exceeding 200\% by the end of the Napoleonic
Wars without default---a feat that Hume and Smith had predicted was
impossible. The key was institutional credibility: Parliament's statutory
control over both taxation and appropriation made British bonds the
safest asset in Europe, driving down borrowing costs even as the
absolute debt grew.

Alexander Hamilton explicitly modeled his 1790 fiscal program on this
British template: a funded debt, a national bank, tariff and excise
revenues, and a sinking fund. Ferguson (1961) describes this as
Hamilton's creation of an American fiscal--military state. But
Hamilton's successors dismantled much of this apparatus: the First Bank
of the United States expired in 1811, the Second Bank was destroyed by
Jackson in 1836, and the federal government operated through the
Independent Treasury system---essentially a government vault with no
banking functions---from 1846 until the Civil War.

When the war came, Chase and Congress had to reconstruct a fiscal--military
state almost from scratch, and the model they reached for---consciously or
not---was once again the British one.

%-----------------------------------------------------------------------
\section{Parallel 1: Suspension of Specie Payments}
\label{sec:suspension}
%-----------------------------------------------------------------------

\subsection{Britain: The Bank Restriction Period, 1797--1821}

On February~26, 1797, facing a drain on gold reserves caused by war
expenditures, French invasion threats, and domestic bank runs, the Privy
Council issued an Order directing the Bank of England to suspend cash
payments. Parliament confirmed this with the Bank Restriction Act of
May~3, 1797. For the next twenty-four years, Bank of England notes were
inconvertible---they could not be exchanged for gold on demand. The notes
continued to circulate, but at a fluctuating discount against gold,
reaching roughly 30\% depreciation by 1813 before recovering as the war
ended.

The Restriction was presented as a temporary emergency measure. Pitt
told Parliament it would last only until ``the dangers which have
produced it are at an end.'' In the event, resumption was repeatedly
postponed---in 1802, 1803, 1814, and 1816---and occurred only in 1821,
six years after Waterloo, at the pre-war gold parity.

\subsection{The Union: Greenback Suspension, 1861--1879}

On December~30, 1861, the major New York banks suspended specie
payments, and the Treasury followed. The Legal Tender Act of
February~25, 1862 then authorized \$150~million in United States Notes
(greenbacks), declared legal tender for all debts ``except duties on
imports and interest as aforesaid.'' Like Bank of England notes after
1797, greenbacks circulated at a fluctuating discount against gold---at
their nadir in July~1864, a greenback dollar was worth only about 39
cents in gold, a depreciation of roughly 61\%.

The parallel is remarkably close:

\begin{center}
\begin{tabular}{lll}
\toprule
\textbf{Feature} & \textbf{Britain 1797} & \textbf{Union 1862} \\
\midrule
Trigger & War-related gold drain & War-related gold drain \\
Legal form & Privy Council order, & Legal Tender Act, \\
 & confirmed by Parliament & passed by Congress \\
Presented as & Temporary emergency & ``A war measure, \\
 & & not of choice'' \\
Peak depreciation & $\sim$30\% (1813) & $\sim$61\% (July 1864) \\
Resumption & 1821 (24 years later) & 1879 (17 years later) \\
Resumption parity & Pre-war gold parity & Pre-war gold parity \\
\bottomrule
\end{tabular}
\end{center}

In both cases, the wartime government faced the same fundamental
trade-off: immediate fiscal necessity against long-run credibility. In
both cases, the government chose suspension, presented it as temporary,
and---crucially---eventually resumed at the pre-war parity, vindicating
bondholders who had lent during the depreciation. The commitment to
resume at par, rather than at the depreciated market rate, is what Bordo
and Kydland (1995) identified as the defining logic of the classical
gold standard as a ``contingent rule'': suspension was permitted during
war emergencies, but resumption at the original parity was the implicit
contract that sustained the government's borrowing capacity.

Henry Adams, writing in 1867---with the Civil War greenback experience
fresh---published ``The Bank of England Restriction, 1797--1821'' (in
\textit{Historical Essays}, 1891), explicitly drawing the parallel. For
Adams, the British Restriction Period was the template through which the
American greenback episode should be understood.

%-----------------------------------------------------------------------
\section{Parallel 2: The Income Tax}
\label{sec:income-tax}
%-----------------------------------------------------------------------

\subsection{Britain: Pitt's Income Tax of 1799}

William Pitt the Younger introduced the world's first modern income tax
in the Act of January~9, 1799 (39~Geo.~III, c.~13), explicitly to
finance the war against revolutionary France. The tax applied a
graduated rate of up to 10\% (two shillings in the pound) on incomes
above \pounds 60, with abatements for incomes between \pounds 60 and
\pounds 200.

The tax was deeply unpopular but enormously productive, raising some
\pounds 6~million in its first year. Addington restructured it in 1803
with withholding at source (Schedule taxation), which became the
template for all subsequent British income taxes. Pitt's tax was repealed
at the Peace of Amiens (1802), reinstated in 1803, and abolished again
after Waterloo in 1816. It was revived by Peel in 1842 and has been
continuous ever since.

\subsection{The Union: Revenue Act of 1861 and Internal Revenue Act of 1862}

The Union introduced America's first income tax with the Revenue Act of
August~5, 1861, which imposed a flat 3\% tax on incomes above \$800.
This was replaced by the Internal Revenue Act of July~1, 1862, which
imposed a graduated structure: 3\% on incomes between \$600 and
\$10{,}000, and 5\% on incomes above \$10{,}000 (later raised to 10\%
on incomes above \$5{,}000 in 1864). The 1862 Act also created the
Bureau of Internal Revenue---a new federal revenue bureaucracy modeled,
in its aspirations, on the professional British revenue services that
Brewer had identified as a hallmark of the fiscal--military state.

The parallels are structural, not merely coincidental:

\begin{enumerate}
\item Both taxes were introduced as emergency wartime measures to
  supplement borrowing capacity.
\item Both employed graduated rates with exemptions for lower incomes.
\item Both were presented as temporary and were repealed after the war
  (the US income tax lapsed in 1872; Britain's in 1816), only to be
  revived later as permanent features of the fiscal system (1842 in
  Britain, 1913 in the United States via the Sixteenth Amendment).
\item Both required the creation of new administrative machinery for
  assessment and collection that outlasted the tax itself.
\end{enumerate}

The intellectual link is suggestive though not fully documented.
American policymakers in 1861 were certainly aware of the British
precedent. Justin Morrill of Vermont, who introduced the first Union
income tax bill, was a voracious reader of British political economy;
the Ways and Means Committee had access to British parliamentary
reports. The British experience demonstrated that a direct tax on
income was both administratively feasible and fiscally productive in
wartime---a proof of concept that made the American innovation less
radical than it might otherwise have appeared.

%-----------------------------------------------------------------------
\section{Parallel 3: Excise Taxes and Internal Revenue}
\label{sec:excise}
%-----------------------------------------------------------------------

\subsection{Britain: The Excise as the ``Hidden Tax''}

The British excise---taxes on domestically produced goods including beer,
spirits, soap, candles, glass, paper, leather, and dozens of other
items---was the workhorse of British war finance from the 1690s onward.
By the Napoleonic Wars, excise revenue roughly equaled customs revenue,
and the two together provided the cash flow that serviced the national
debt. Brewer emphasized that the excise was uniquely effective because
it was collected at the point of manufacture by a professional
bureaucracy, making evasion difficult.

\subsection{The Union: The Internal Revenue Act of 1862}

The Union's Internal Revenue Act of 1862 was the most sweeping tax
legislation in American history to that date. Beyond the income tax, it
imposed excise duties on virtually every manufactured article: spirits,
tobacco, beer, patent medicines, playing cards, billiard tables,
carriages, yachts, slaughtered cattle, manufactured cotton, iron, steel,
paper, leather, and dozens more. Stamp duties were applied to legal
documents, bank checks, insurance policies, telegrams, and photographs.
License fees were imposed on bankers, brokers, pawnbrokers, distillers,
and auctioneers.

The breadth of the 1862 Act echoed the British excise system in its
ambition to tax every physical manifestation of economic activity. David
Wells, appointed by Lincoln as Special Commissioner of the Revenue in
1866, explicitly compared the American system to the British excise and
recommended rationalizing it along British lines---reducing the number
of taxed articles and concentrating on a few high-yield items (spirits,
tobacco, beer). This is precisely what happened during Reconstruction:
by the mid-1870s, the internal revenue system had been simplified to
rely primarily on taxes on whiskey, tobacco, and beer---the same three
goods that dominated British excise collections.

%-----------------------------------------------------------------------
\section{Parallel 4: Long-Term Funded Debt and the Sinking Fund}
\label{sec:funded-debt}
%-----------------------------------------------------------------------

\subsection{Britain: Consols and Pitt's Sinking Fund}

The cornerstone of British public credit was the consolidated annuity
(``consol''), created in 1751 by consolidating several older
obligations into a single perpetual bond paying 3\% interest. Consols
had no maturity date---they were perpetuities, redeemable at the
government's option but never obligatorily. Their liquidity and
uniformity made them the benchmark asset of the London capital market
and, by extension, of the global financial system.

To manage the accumulation of war debt, Pitt established a new Sinking
Fund in 1786, building on earlier models by Walpole (1717) and Pelham
(1750). The Fund received a fixed annual appropriation from tax
revenues, which was used to purchase government bonds on the open market
for cancellation. The logic was that dedicated revenues, protected from
parliamentary diversion, would provide a credible commitment to
gradual debt retirement, sustaining investor confidence even during
wartime borrowing. In practice, the Sinking Fund was undermined during
the Napoleonic Wars when the government simultaneously borrowed to fund
the Sinking Fund's purchases---what Ricardo denounced as an accounting
fiction. Nevertheless, the \textit{existence} of the Fund served a
signaling function: it communicated the government's intention to
retire debt after the war.

\subsection{The Union: 5-20 Bonds and the Statutory Sinking Fund}

The Union's 5-20 bonds---authorized by Section~2 of the Legal Tender
Act of February~25, 1862---were not perpetuities like consols, but they
shared key structural features. They bore 6\% interest, were redeemable
at the government's pleasure after five years, and payable in twenty
years. The callable feature gave the Treasury the same optionality that
the British government enjoyed with consols: the right to refinance when
interest rates fell, without being locked into expensive coupon
payments.

Section~5 of the Legal Tender Act established a statutory sinking
fund strikingly reminiscent of Pitt's. Customs duties---required to be
paid in coin---were dedicated first to the payment of interest on bonds
and notes in coin, and second to a sinking fund of 1\% of the total
outstanding debt per year for gradual principal retirement. The
earmarking of a specific revenue source (customs) for a specific purpose
(debt service and retirement) replicated the logic of the British
system, in which specific taxes were ``hypothecated'' to specific debt
obligations.

The sinking fund served the same dual function in both episodes: it
provided a genuine (if modest) mechanism for debt reduction, and---more
importantly---it signaled fiscal credibility to investors. As documented
in the companion note on the 5-20 bonds, Jay Cooke's bond-selling
agents used the sinking fund clause to assure investors that the bonds
would be redeemed in coin, even though the statute was silent on the
currency of principal repayment. The sinking fund, like its British
antecedent, was as much a commitment device as a fiscal mechanism.

%-----------------------------------------------------------------------
\section{Parallel 5: The Central Bank as Fiscal Agent}
\label{sec:central-bank}
%-----------------------------------------------------------------------

\subsection{Britain: The Bank of England}

The Bank of England was founded in 1694 explicitly to lend to the
government: its original charter provided a \pounds 1.2~million loan to
William III in exchange for a banking monopoly. By the Napoleonic Wars,
the Bank served as the government's fiscal agent (managing bond
issuance and debt service), its banker (holding government deposits and
transferring funds), and its monetary authority (controlling the note
issue). The symbiosis was intimate: the government depended on the Bank
for war finance, and the Bank depended on its government relationship
for its charter and its monopoly.

During the Restriction Period, the Bank lent massively to the
government by purchasing Exchequer bills and discounting Treasury
paper---in effect, monetizing the war deficit. This provoked the
Bullionist Controversy of 1810, in which Ricardo and others argued that
the Bank's note issue, freed from the discipline of gold convertibility,
had depreciated the currency beyond what military necessity required.

\subsection{The Union: The National Banking System as Surrogate}

The Union had no central bank in 1861. The Second Bank of the United
States had been destroyed in 1836, and the Independent Treasury system
(1846--1913) deliberately separated government finance from banking.
Chase was acutely aware of this institutional gap.

The National Banking Acts of 1863--1864 were Chase's partial solution.
While they did not create a central bank, they created a network of
nationally chartered banks required to hold U.S.\ government bonds as
backing for their note issues. This accomplished two objectives
simultaneously: it created a captive market for government bonds
(increasing demand and supporting prices) and it created a uniform
national currency (national bank notes) to replace the chaotic
state-bank note system. The 10\% tax on state bank notes (imposed in
1865) completed the nationalization of the currency.

The structural logic mirrors the Bank of England relationship, with the
roles distributed across an entire banking system rather than
concentrated in a single institution:

\begin{center}
\begin{tabular}{p{5cm}p{5cm}}
\toprule
\textbf{Bank of England} & \textbf{National Banking System} \\
\midrule
Lends to government by & Required to hold government \\
purchasing Exchequer bills & bonds as note-backing \\[0.5em]
Creates money backed by & Creates national bank notes \\
government obligations & backed by government bonds \\[0.5em]
Government deposits held & Government deposits held \\
at the Bank & at designated national banks \\[0.5em]
Single institution with & Distributed system with \\
monopoly charter & federal charters \\
\bottomrule
\end{tabular}
\end{center}

The key difference is that the British model concentrated fiscal agency,
note issuance, and monetary control in a single institution, whereas the
American model disaggregated these functions---a design choice that
reflected Jacksonian hostility to centralized financial power but
created the ``inelasticity'' problem that would plague the national
banking system until the creation of the Federal Reserve in 1913.

%-----------------------------------------------------------------------
\section{Parallel 6: Customs Duties as the Fiscal Anchor}
\label{sec:customs}
%-----------------------------------------------------------------------

\subsection{Britain: Customs as Debt-Service Revenue}

British customs duties served a dual function throughout the 18th and
early 19th centuries: they were a major source of general revenue, and
they were specifically pledged to service particular debt obligations.
The practice of ``hypothecation''---dedicating a specific tax to a
specific spending purpose---meant that investors could assess the
security of a bond by evaluating the reliability of the tax stream
backing it. Customs revenues, collected on a large volume of international
trade funneled through a manageable number of ports, were among the most
predictable and hardest to evade of all revenue sources.

\subsection{The Union: Coin Duties as the Gold Anchor}

Section~5 of the Legal Tender Act required import duties to be paid in
coin---not in depreciated greenbacks---and earmarked this coin revenue
for interest payments on government bonds and the sinking fund. This
provision created what Lowenstein (2022) called an ``alchemy machine'':
it guaranteed that bondholders' coupon income would be paid in gold
regardless of the depreciation of the paper currency. The customs-coin
requirement was the hard anchor of the entire Union credit system.

The structural logic is identical to the British practice of
hypothecation. Just as British investors could evaluate a bond by
examining the specific tax revenue backing it, American investors could
evaluate the 5-20 bonds by examining the customs revenue pledged to
their interest payments. The coin requirement added an additional
dimension absent from the British case: not only was a specific revenue
stream dedicated to debt service, but it was denominated in a different
(harder) currency than the medium of general circulation.

This is the feature that created the asymmetry at the heart of the 5-20
controversy: statutory commitment for interest (guaranteed in coin by
Section~5) versus merely reputational commitment for principal
(implied by Treasury marketing but absent from the statute). The British
system, by contrast, operated in a single currency (Bank of England
notes), with the commitment running to eventual resumption of
convertibility rather than to a dual currency structure within the
system.

%-----------------------------------------------------------------------
\section{Parallel 7: Post-War Resumption and the Commitment to Creditors}
\label{sec:resumption}
%-----------------------------------------------------------------------

\subsection{Britain: Resumption at Par, 1819--1821}

After Waterloo, the British government faced the question that every
wartime suspender must eventually confront: at what rate to resume
convertibility. Resumption at the pre-war parity (the old mint price of
\pounds 3\,17s\,10\textonehalf d per troy ounce of gold) would reward
creditors who had purchased bonds during the depreciation, but it would
impose deflation on debtors and the general economy. Resumption at the
depreciated market rate would ease the adjustment but would constitute a
partial repudiation of the government's implicit promise to restore the
old standard.

The Bullion Committee of 1810, Ricardo's \textit{Proposals for an
Economical and Secure Currency} (1816), and ultimately Peel's Act of
1819 all pointed in the same direction: resumption at the old parity.
This was accomplished gradually between 1819 and 1821. The deflationary
consequences were severe---agricultural distress, manufacturing
contraction, political unrest---but the government's credit was
vindicated. British consols, which had traded well below par during the
war, recovered to par and above.

\subsection{The Union: Resumption at Par, 1875--1879}

The Union faced an identical choice. By 1865 a greenback dollar was
worth roughly 70 cents in gold. Pendleton's Ohio Idea and the 1868
Democratic platform argued for paying bond principal in greenbacks
(effectively, resumption at the depreciated rate). Grant and the
Republicans insisted on payment in coin---and the Public Credit Act of
1869 committed the government's ``solemn'' faith to coin payment.

The Resumption Act of January~14, 1875 set January~1, 1879 as the date
for full resumption of specie payments on greenbacks at par. Treasury
Secretary John Sherman built a gold reserve to back the resumption
commitment. When the date arrived, virtually no one actually presented
greenbacks for redemption---the credibility of the commitment was
sufficient to maintain parity without significant gold outflows.

The logic is identical to the British case: the government chose
deflation and creditor vindication over debtor relief and partial
repudiation. In both cases, the post-war commitment debate was
ultimately a referendum on whether the government's wartime borrowing
should be treated as state-contingent (with the currency of repayment
adjusting to war outcomes) or as an absolute obligation (with the
pre-war standard restored regardless of interim depreciation).

Bordo and Kydland (1995) formalized this as the ``contingent rule''
interpretation of the classical gold standard: the gold standard was
a rule with an escape clause for wartime emergencies, credible precisely
because governments consistently resumed at the pre-war parity after
the emergency passed. Both Britain after 1815 and the United States
after 1865 were key data points for this interpretation.

%-----------------------------------------------------------------------
\section{Channels of Transmission: Conscious Emulation or Convergent Evolution?}
%-----------------------------------------------------------------------

How much of the British--Union parallelism reflects deliberate imitation,
and how much reflects similar institutional responses to similar fiscal
pressures? The evidence suggests a mixture.

\subsection{Conscious Emulation}

Several channels of direct influence are identifiable:

\begin{enumerate}
\item \textbf{Hamilton's precedent.} The entire Hamiltonian fiscal
  system---the template that Civil War policymakers were reconstructing
  under pressure---was itself explicitly modeled on British practice.
  Chase, Spaulding, and the Ways and Means Committee were therefore
  drawing on a British model at one remove: they were restoring
  Hamilton's version of Pitt's fiscal state.

\item \textbf{Intellectual currency of British political economy.}
  American statesmen of the 1860s read Smith, Ricardo, McCulloch, and
  Mill. The Bullionist Controversy, the Restriction Period, and Peel's
  resumption were well-known episodes. When Fessenden said the Legal
  Tender Act ``shocks all my notions of political, moral, and national
  honor,'' he was speaking in a register shaped by the British classical
  economists' horror of inconvertible paper.

\item \textbf{The British income tax as proof of concept.} Pitt's income
  tax demonstrated that a direct tax on income was administratively
  feasible and fiscally productive in a common-law country during an
  existential war. Without this precedent, it is unclear whether
  Congress would have adopted the form so readily in 1861--1862.

\item \textbf{Henry Adams's explicit comparison.} Adams's 1867 essay on
  ``The Bank of England Restriction'' was written precisely to draw
  lessons from the British experience for the American greenback
  question. Adams was not an obscure figure: he was the grandson and
  great-grandson of presidents, writing in the \textit{North American
  Review}. His essay demonstrates that the comparison was live in
  contemporaneous policy discussion.
\end{enumerate}

\subsection{Convergent Evolution}

Other parallels likely reflect structural necessity rather than
conscious imitation:

\begin{enumerate}
\item \textbf{Suspension} is the natural response when a government
  committed to metallic money faces wartime expenditures that exhaust
  its gold reserves. Any gold-standard country fighting an expensive war
  will arrive at this point; conscious emulation is unnecessary.

\item \textbf{Excise taxation} is the natural domestic revenue source
  when customs duties are already committed to debt service and an
  income tax is politically constrained. The Union arrived at a
  comprehensive excise system not because it was copying the British
  excise, but because the logic of wartime revenue maximization pointed
  in the same direction.

\item \textbf{Post-war resumption at par} reflects a logic intrinsic to
  sovereign borrowing under a metallic standard: the government that
  does not resume at par loses access to cheap credit in the next war.
  The ``contingent rule'' operates whether or not policymakers are
  consciously aware of the British precedent.
\end{enumerate}

The most likely characterization is that the Union's fiscal program was
shaped by British precedents \textit{filtered through Hamilton}, with
the most direct imitation occurring in the institutional designs
(national banking, sinking fund, income tax structure) and the most
purely structural convergence occurring in the monetary responses
(suspension, depreciation, resumption).

%-----------------------------------------------------------------------
\section{Analytical Implications}
%-----------------------------------------------------------------------

The British--Union parallelism enriches several analytical themes
developed in the companion notes on the 5-20 bonds.

\subsection{Commitment Theory}

The core insight of both episodes is that wartime governments face a
commitment problem: they need to borrow at non-ruinous rates, but
wartime pressures undermine the credibility of their repayment promises.
Both Britain and the Union addressed this problem through a
\textit{layered commitment structure}:

\begin{itemize}
\item \textbf{Statutory commitment:} Specific tax revenues dedicated to
  specific debt-service obligations (British hypothecation; the Union's
  coin-customs/interest linkage in Section~5).
\item \textbf{Institutional commitment:} A banking system whose own
  solvency depended on government creditworthiness (the Bank of England;
  national banks required to hold government bonds).
\item \textbf{Reputational commitment:} Extra-legal assurances by
  Treasury officials (Chase's gold-payment marketing of the 5-20s;
  Pitt's repeated promises that the Restriction was temporary).
\item \textbf{Post-war vindication:} Resumption at par, confirming the
  implicit contract and sustaining credibility for future borrowing.
\end{itemize}

This layered structure is precisely the Kydland--Prescott (1977) and
Rogoff (1985) problem of credible commitment, solved not through a
single mechanism but through a mutually reinforcing combination of
statutory, institutional, and reputational devices.

\subsection{State-Contingent Debt}

Both episodes created something close to the Lucas and Stokey (1983)
state-contingent debt structure, despite operating within legal and
ideological frameworks that formally rejected contingency. British
bondholders during the Restriction Period received interest payments in
depreciated Bank of England notes---effectively, the real value of their
coupons was contingent on the war's progress. Union bondholders' coupon
income was protected by the coin clause, but the real value of their
principal was contingent on two uncertainties: the currency of repayment
(coin vs.\ greenbacks) and the greenback--gold exchange rate. In both
cases, the debt was \textit{de facto} state-contingent even as it was
\textit{de jure} fixed-nominal.

The difference is instructive. Because Britain operated in a single
currency (Bank of England notes), the contingency operated through the
real value of the currency itself---a form of implicit default via
inflation. Because the Union operated in a dual-currency system (coin
for customs and interest; greenbacks for everything else), the
contingency operated through the \textit{exchange rate} between the two
currencies---a more visible and politically explosive form of the same
phenomenon.

\subsection{Bryant--Wallace Price Discrimination}

Both governments practiced fiscal discrimination among creditor classes,
in the spirit of Bryant and Wallace (1984). British creditors during the
Restriction received depreciated-currency coupons but expected par
resumption; British taxpayers bore the burden of deflation during
resumption. Union soldiers were paid in greenbacks; bondholders received
gold coupons; domestic small investors held bonds with ambiguous principal
commitments; foreign investors received gold throughout by convention.
In both cases, the fiscal burden of war was distributed unequally across
classes, with the distribution mediated by the institutional structure
of the currency and debt system.

\subsection{The Rogoff Conservative Central Banker}

Lincoln's appointment of Chase---a hard-money, former Jackson Democrat
whose entire political identity was bound up with sound finance---as
the Treasury Secretary who would oversee legal tender issues and
inconvertible paper money parallels the Rogoff (1985) logic of
delegating monetary policy to a ``conservative'' central banker. Chase's
personal revulsion at paper money (\textit{exactly} paralleling
Fessenden's famous denunciation in the Senate) served as a commitment
device: investors could believe that a man who hated inconvertible paper
would work to restore convertibility as soon as the emergency passed.
Pitt performed the same function in Britain: his reputation for fiscal
rectitude made the Restriction more palatable to investors, precisely
because Pitt was known to regard it as a regrettable necessity.

%-----------------------------------------------------------------------
\section{A Comparative Table}
%-----------------------------------------------------------------------

\begin{center}
\small
\begin{longtable}{p{3.8cm}p{4.5cm}p{4.5cm}}
\toprule
\textbf{Element} & \textbf{Britain, 1793--1821} & \textbf{Union, 1861--1879} \\
\midrule
\endfirsthead
\toprule
\textbf{Element} & \textbf{Britain, 1793--1821} & \textbf{Union, 1861--1879} \\
\midrule
\endhead
\bottomrule
\endlastfoot
Specie suspension &
  Bank Restriction Act, 1797; resumed 1821 at pre-war parity &
  Suspension Dec.~1861; Legal Tender Act Feb.~1862; resumed Jan.~1879 at pre-war parity \\[0.5em]
Inconvertible paper &
  Bank of England notes, legal tender for all purposes &
  Greenbacks, legal tender except customs duties and bond interest \\[0.5em]
Peak depreciation &
  $\sim$30\% (1813) &
  $\sim$61\% (July 1864) \\[0.5em]
Income tax &
  Pitt's tax, 1799 (up to 10\%); repealed 1816; restored 1842 &
  Revenue Act 1861 (3\% flat); graduated 1862--1864 (3--10\%); lapsed 1872; restored 1913 \\[0.5em]
Comprehensive excise &
  Extensive excise system on domestic manufactures; professional Excise service &
  Internal Revenue Act 1862; Bureau of Internal Revenue created; simplified post-war \\[0.5em]
Customs earmarking &
  Customs revenues hypothecated to specific debt obligations &
  Customs in coin earmarked to bond interest and sinking fund (Section~5) \\[0.5em]
Funded debt instrument &
  Consols: 3\% perpetuities, callable at government option &
  5-20 bonds: 6\% bonds, callable after 5 years, maturing in 20 \\[0.5em]
Sinking fund &
  Pitt's Sinking Fund (1786): fixed annual purchase of bonds from earmarked revenue &
  Section~5 sinking fund: 1\% of total debt per year from coin customs revenue \\[0.5em]
Central bank / banking &
  Bank of England: fiscal agent, government banker, note issuer &
  National Banking System: bond-backed notes, government depositories, no single institution \\[0.5em]
Post-war commitment &
  Peel's Act (1819): resumption at old parity by 1821 &
  Public Credit Act (1869) + Resumption Act (1875): resumption at par by Jan.~1, 1879 \\[0.5em]
Post-war politics &
  Agricultural distress; anti-resumption agitation; Radical reform movement &
  Ohio Idea; Greenback Party; Bland--Allison Act; populist agitation \\[0.5em]
\end{longtable}
\end{center}

%-----------------------------------------------------------------------
\section{Conclusion: Pitt's Ghost in Lincoln's Treasury}
%-----------------------------------------------------------------------

The Union's Civil War fiscal program was, in its essential architecture,
a recapitulation of Britain's Napoleonic War fiscal program---adapted to
American institutional conditions and constrained by Jacksonian
suspicion of centralized financial power. Where Pitt had the Bank of
England, Chase created a distributed national banking system. Where
Pitt had consols, Chase had 5-20 bonds. Where Pitt had a
comprehensive excise, Chase had the Internal Revenue Act. Where Pitt
had the income tax, Chase had the Revenue Acts of 1861 and 1862.
Where Pitt had the Sinking Fund, Chase had Section~5.

In every case, the British precedent established the template---either
directly, through Hamilton's conscious imitation a generation earlier;
indirectly, through the intellectual influence of British political
economy on American statesmen; or structurally, through the logic of
wartime fiscal necessity under a metallic monetary standard. The most
important shared feature was not any single policy instrument but the
\textit{commitment architecture}: the layered combination of statutory
dedication of revenues, institutional symbiosis between banks and
government bonds, reputational assurances by credible Treasury
officials, and post-war resumption at pre-war parity.

Understanding the British precedent illuminates the 5-20 bond episode
in particular. The ambiguity over coin versus greenback repayment of
principal---the central puzzle of the companion document---was itself a
variant of the ambiguity that pervaded British war finance during the
Restriction Period: would the government honor its implicit promise to
return to gold at the old parity, or would it settle for a depreciated
resumption? In both cases, the ambiguity served a political function
(preserving the wartime coalition) and a fiscal function (maintaining
borrowing capacity under uncertainty). In both cases, post-war politics
eventually resolved the ambiguity in favor of creditors---at
significant cost to debtors and taxpayers.

The British precedent thus provides both a comparative case and a causal
antecedent for the Union's fiscal design. It suggests that the 5-20
bond structure, the Legal Tender Act, the National Banking Acts, and the
post-war resumption program should be understood not as ad hoc responses
to crisis but as elements of a coherent fiscal--military state
template, developed in Britain over a century of warfare and imported
into American practice through multiple channels of intellectual and
institutional transmission.

%-----------------------------------------------------------------------
\section*{References}
%-----------------------------------------------------------------------

\begin{description}
\setlength\itemsep{0.5em}

\item Adams, Henry. 1891 [1867]. ``The Bank of England Restriction,
1797--1821.'' In \textit{Historical Essays}, pp.~178--236. New York:
Charles Scribner's Sons.

\item Bordo, Michael D. and Finn E. Kydland. 1995. ``The Gold Standard
As a Rule: An Essay in Exploration.'' \textit{Explorations in Economic
History} 32(4): 423--464.

\item Brewer, John. 1989. \textit{The Sinews of Power: War, Money,
and the English State, 1688--1783}. New York: Alfred A.~Knopf.

\item Bryant, John and Neil Wallace. 1984. ``A Price Discrimination
Analysis of Monetary Policy.'' \textit{Review of Economic Studies}
51(2): 279--288.

\item Daunton, Martin. 2001. \textit{Trusting Leviathan: The Politics
of Taxation in Britain, 1799--1914}. Cambridge: Cambridge University
Press.

\item Dickson, P.~G.~M. 1967. \textit{The Financial Revolution in England:
A Study in the Development of Public Credit, 1688--1756}. London:
Macmillan.

\item Fetter, Frank W. 1965. \textit{Development of British Monetary
Orthodoxy, 1797--1875}. Cambridge, MA: Harvard University Press.

\item Ferguson, E.~James. 1961. \textit{The Power of the Purse: A
History of American Public Finance, 1776--1790}. Chapel Hill:
University of North Carolina Press.

\item Hall, George J. and Thomas J. Sargent. 2014. ``Fiscal
Discriminations in Three Wars.'' \textit{Journal of Monetary Economics}
61: 148--166.

\item Hall, George J. and Thomas J. Sargent. 2022. ``Three World Wars:
Fiscal--Monetary Consequences.'' \textit{Proceedings of the National
Academy of Sciences} 119(18): e2200349119.

\item Kydland, Finn E. and Edward C. Prescott. 1977. ``Rules Rather
Than Discretion: The Inconsistency of Optimal Plans.'' \textit{Journal
of Political Economy} 85(3): 473--492.

\item Lowenstein, Roger. 2022. \textit{Ways and Means: Lincoln and His
Cabinet and the Financing of the Civil War}. New York: Penguin Press.

\item Lucas, Robert E., Jr. and Nancy L. Stokey. 1983. ``Optimal Fiscal
and Monetary Policy in an Economy Without Capital.'' \textit{Journal of
Monetary Economics} 12(1): 55--93.

\item Mitchell, Wesley C. 1903. \textit{A History of the Greenbacks}.
Chicago: University of Chicago Press.

\item O'Brien, Patrick K. 1988. ``The Political Economy of British
Taxation, 1660--1815.'' \textit{Economic History Review} 41(1): 1--32.

\item O'Brien, Patrick K. and Philip A. Hunt. 1993. ``The Rise of a
Fiscal State in England, 1485--1815.'' \textit{Historical Research}
66(160): 129--176.

\item Ricardo, David. 1816. \textit{Proposals for an Economical and
Secure Currency; With Observations on the Profits of the Bank of
England}. London: John Murray.

\item Rockoff, Hugh. 2012. \textit{America's Economic Way of War: War
and the US Economy from the Spanish-American War to the Persian Gulf
War}. Cambridge: Cambridge University Press.

\item Rogoff, Kenneth. 1985. ``The Optimal Degree of Commitment to an
Intermediate Monetary Target.'' \textit{Quarterly Journal of Economics}
100(4): 1169--1189.

\item Sargent, Thomas J. and Fran\c{c}ois R. Velde. 1995. ``Macroeconomic
Features of the French Revolution.'' \textit{Journal of Political
Economy} 103(3): 474--518.

\item Studenski, Paul and Herman E. Krooss. 1952. \textit{Financial
History of the United States}. New York: McGraw-Hill.

\item Unger, Irwin. 1964. \textit{The Greenback Era: A Social and
Political History of American Finance, 1865--1879}. Princeton:
Princeton University Press.

\end{description}

\end{document}
